\KOMAoption{paper}{a4, portrait}
\areaset{180mm}{267mm}
\fancyheadoffset{0pt} % обновляем хедер и футер

\newgeometry{
  a4paper,
  top=15mm,
  bottom=10mm,
  right=5mm,
  left=20mm,
}

\color{uniblue}{\section[(обязательное) Карта заказа]{(обязательное)\\Карта заказа}}
\color{black}
\scriptsize
\hspace{6cm} % Добавляем отступ слева в 2 см
\textbf{ЮНИТ-М319-ДЗТ2-
    \tikz[baseline=(char.base)]{
        \node[rectangle, fill=uniblue, text=white, inner sep=2pt] (char) {\textbf{M1}};
    } -
    \tikz[baseline=(char.base)]{
        \node[rectangle, fill=unired, text=white, inner sep=2pt] (char) {\textbf{M2}};
    } -
    \tikz[baseline=(char.base)]{
        \node[rectangle, fill=uniblue, text=white, inner sep=2pt] (char) {\textbf{M3}};
    } - 
    \tikz[baseline=(char.base)]{
        \node[rectangle, fill=uniblue, text=white, inner sep=2pt] (char) {\textbf{M4}};
    } - 
    \tikz[baseline=(char.base)]{
        \node[rectangle, fill=uniblue, text=white, inner sep=2pt] (char) {\textbf{M5}};
    } - 
    \tikz[baseline=(char.base)]{
        \node[rectangle, fill=uniblue, text=white, inner sep=2pt] (char) {\textbf{M6}};
    } -  
    \tikz[baseline=(char.base)]{
        \node[rectangle, fill=uniblue, text=white, inner sep=2pt] (char) {\textbf{M7}};
    } - 
    \tikz[baseline=(char.base)]{
        \node[rectangle, fill=unired, text=white, inner sep=2pt] (char) {\textbf{M8}};
    } -
    \tikz[baseline=(char.base)]{
        \node[rectangle, fill=uniblue, text=white, inner sep=2pt] (char) {\textbf{M9}};
    } - 
    \tikz[baseline=(char.base)]{
        \node[rectangle, fill=unired, text=white, inner sep=2pt] (char) {\textbf{M10}};
    } - 
    \tikz[baseline=(char.base)]{
        \node[rectangle, fill=unired, text=white, inner sep=2pt] (char) {\textbf{Ф}};
    } - 
    \tikz[baseline=(char.base)]{
        \node[rectangle, fill=uniblue, text=white, inner sep=2pt] (char) {\textbf{В}};
    }                                     
}

\scriptsize
\vspace{0.3cm}

\begin{flushleft}
    \begin{tabular}{|p{6cm}|p{2.1cm}|}
        \hline
        \rowcolor{uniblue} % Синий цвет с прозрачностью
        \multicolumn{2}{|l|}{\textcolor{white}{Модуль в слоте М1}} \\ % Объединение ячеек
        \hline
        Модуль источника питания Uном=110 В & Р01 \\ 
        \hline
        Модуль источника питания Uном=220 В    & Р02 \\
        \hline
        Модуль источника питания Uном=220 В с блоком конденсаторов    & Р02c \\
        \hline
        Модуль источника питания Uном=24 В    & Р04 \\
        \hline
    \end{tabular}
\end{flushleft}

% Начинаем TikZ
\begin{tikzpicture}[overlay]
    \draw[uniblue, very thick] (8.45, 2.7) -- (8.45, 2.05) -- (7.97, 2.05); % Новые координаты
\end{tikzpicture}

\begin{flushleft}
    \begin{tabular}{|p{6cm}|p{2.1cm}|}
        \hline
        \rowcolor{unired}
        \multicolumn{2}{|l|}{\textcolor{white}{Модуль в слоте М2}} \\ % Объединение ячеек
        \hline
        Нет модуля или в слоте установлены модули: P05, P02c & X \\ 
        \hline
        Модуль дискретных входов    & B001,B002,B003,B021 \\
        \hline
        Модуль реле    & K001,K002,K003,K004 \\
        \hline
    \end{tabular}
\end{flushleft}

% Начинаем TikZ
\begin{tikzpicture}[overlay]
    \draw[unired, very thick] (9.13, 4.4) -- (9.13, 1.35) -- (7.97, 1.35); % Новые координаты
\end{tikzpicture}


\begin{flushleft}
    \begin{tabular}{|p{6cm}|p{2.1cm}|}
        \hline
        \rowcolor{uniblue} % Синий цвет с прозрачностью
        \multicolumn{2}{|l|}{\textcolor{white}{Модуль в слоте М3...M7}} \\ % Объединение ячеек
        \hline
        Нет модуля & Х \\ 
        \hline
        Модуль дискретных входов    & B001,B002,B003,B021 \\
        \hline
        Модуль реле    & K001,K002,K003,K004 \\
        \hline
    \end{tabular}
\end{flushleft}

% Начинаем TikZ
\begin{tikzpicture}[overlay]
    \draw[uniblue, very thick] (9.8, 6.05) -- (9.8, 1.35) -- (7.97, 1.35); % Новые координаты
\end{tikzpicture}


\begin{flushleft}
    \begin{tabular}{|p{6cm}|p{2.1cm}|}
        \hline
        \rowcolor{unired}
        \multicolumn{2}{|l|}{\textcolor{white}{Модуль в слоте М8}} \\ % Объединение ячеек
        \hline
        Нет модуля или в слоте установлены модули: P05, P02c & X \\ 
        \hline
        Модуль дискретных входов    & B001,B002,B003,B021 \\
        \hline
        Модуль реле    & K001,K002,K003,K004 \\
        \hline
    \end{tabular}
\end{flushleft}

% Начинаем TikZ
\begin{tikzpicture}[overlay]
    \draw[unired, very thick] (13.13, 7.7) -- (13.13, 1.35) -- (7.97, 1.35); % Новые координаты
\end{tikzpicture}


\begin{flushleft}
    \begin{tabular}{|p{6cm}|p{2.1cm}|}
        \hline
        \rowcolor{uniblue} % Синий цвет с прозрачностью
        \multicolumn{2}{|l|}{\textcolor{white}{Модуль в слоте М9}} \\ % Объединение ячеек
        \hline
        Нет модуля & Х \\ 
        \hline
        Модуль дискретных входов    & B001,B002,B003,B021 \\
        \hline
        Модуль реле    & K001,K002,K003,K004 \\
        \hline
    \end{tabular}
\end{flushleft}

% Начинаем TikZ
\begin{tikzpicture}[overlay]
    \draw[uniblue, very thick] (13.76, 9.4) -- (13.76, 1.35) -- (7.97, 1.35); % Новые координаты
\end{tikzpicture}











\clearpage %!!!!!!!!!!БЕЗ ЭТОГО ХРЕНЬ

\restoregeometry



% КООРДИНАТЫ ДЛЯ ЛИНИЙ
% (А, В) -- (А, С) --(D, С)
%
%                      | B(^)
%                      | 
%    D(<>) ------------A(<>)C(^)

\def\aDelta{4} %смещение координаты А
\def\dCoord{11} % начальная левая координата для всех линий определяется как 7+\aDelta


\newgeometry{
  a4paper,
  top=15mm,
  bottom=10mm,
  right=5mm,
  left=20mm,
}
\begin{landscape}

\setlength{\headsep}{7.5mm} 
\setlength{\topmargin}{-27mm}
\setlength{\footskip}{0mm}

\color{uniblue}{\section[(обязательное) Карта заказа]{(обязательное)\\Карта заказа}}
\color{black}
\scriptsize
\hspace{\dimexpr 6cm + \aDelta cm} % Добавляем отступ слева в 2 см
\textbf{ЮНИТ-М319-ДЗТ2-
    \tikz[baseline=(char.base)]{
        \node[rectangle, fill=uniblue, text=white, inner sep=2pt] (char) {\textbf{M1}};
    } -
    \tikz[baseline=(char.base)]{
        \node[rectangle, fill=unired, text=white, inner sep=2pt] (char) {\textbf{M2}};
    } -
    {\textbf{K002}
    } - 
    {\textbf{K002}
    } - 
    \tikz[baseline=(char.base)]{
        \node[rectangle, fill=uniblue, text=white, inner sep=2pt] (char) {\textbf{M5}};
    } - 
    \tikz[baseline=(char.base)]{
        \node[rectangle, fill=uniblue, text=white, inner sep=2pt] (char) {\textbf{M6}};
    } -  
    \tikz[baseline=(char.base)]{
        \node[rectangle, fill=uniblue, text=white, inner sep=2pt] (char) {\textbf{M7}};
    } - 
    {\textbf{B021}
    } -
    {\textbf{B021}
    } -
    {\textbf{B021}
    } - 
    {\textbf{B001}
    } -  
    \tikz[baseline=(char.base)]{
        \node[rectangle, fill=unired, text=white, inner sep=2pt] (char) {\textbf{M12}};
    } - 
    {\textbf{X}
    } -      
    \tikz[baseline=(char.base)]{
        \node[rectangle, fill=uniblue, text=white, inner sep=2pt] (char) {\textbf{M14}};
    } - 
    \tikz[baseline=(char.base)]{
        \node[rectangle, fill=unired, text=white, inner sep=2pt] (char) {\textbf{Ф}};
    } - 
    \tikz[baseline=(char.base)]{
        \node[rectangle, fill=uniblue, text=white, inner sep=2pt] (char) {\textbf{В}};
    }                                     
}

\scriptsize
\vspace{0.3cm}

\begin{flushleft}
    \hspace{\aDelta cm}
    \begin{tabular}{|p{6.2cm}|p{2.1cm}|}
        \hline
        \rowcolor{uniblue} % Синий цвет с прозрачностью
        \multicolumn{2}{|l|}{\textcolor{white}{Модуль в слоте М1}} \\ % Объединение ячеек
        \hline
        Модуль источника питания Uном=/\~{} 110 В & Р01 \\ 
        \hline
        Модуль источника питания Uном=/\~{} 220 В    & Р02 \\
        \hline
        Модуль источника питания Uном=/\~{} 220 В, блок конденсаторов    & Р02c \\
        \hline
        Модуль источника питания Uном=24 В    & Р04 \\
        \hline
    \end{tabular}
\end{flushleft}

% Начинаем TikZ
\begin{tikzpicture}[overlay]
    \draw[uniblue, very thick] (8.45+\aDelta, 2.7) -- (8.45+\aDelta, 2.05) -- (\dCoord, 2.05); % Новые координаты
\end{tikzpicture}

\begin{flushleft}
    \hspace{\aDelta cm}
    \begin{tabular}{|p{6.2cm}|p{2.1cm}|}
        \hline
        \rowcolor{unired}
        \multicolumn{2}{|l|}{\textcolor{white}{Модуль в слоте М2}} \\ % Объединение ячеек
        \hline
        Нет модуля или в слоте М1 установлены модули: P05, P02c & X \\ 
        \hline
        Модуль дискретных входов    & B001,B002,B003,B021 \\
        \hline
        Модуль реле    & K001,K002,K003,K004 \\
        \hline
    \end{tabular}
\end{flushleft}

% Начинаем TikZ
\begin{tikzpicture}[overlay]
    \draw[unired, very thick] (9.13+\aDelta, 4.4) -- (9.13+\aDelta, 1.35) -- (\dCoord, 1.35); % Новые координаты
\end{tikzpicture}


\begin{flushleft}
    \hspace{\aDelta cm}
    \begin{tabular}{|p{6.2cm}|p{2.1cm}|}
        \hline
        \rowcolor{uniblue} % Синий цвет с прозрачностью
        \multicolumn{2}{|l|}{\textcolor{white}{Модуль в слоте М5...M7}} \\ % Объединение ячеек
        \hline
        Нет модуля & Х \\ 
        \hline
        Модуль дискретных входов    & B001,B002,B003,B021 \\
        \hline
        Модуль реле    & K001,K002,K003,K004 \\
        \hline
    \end{tabular}
\end{flushleft}

% Начинаем TikZ
\begin{tikzpicture}[overlay]
    \draw[uniblue, very thick] (11.4+\aDelta, 6.05) -- (11.4+\aDelta, 1.35) -- (\dCoord, 1.35); 
    \draw[uniblue, very thick] (12.1+\aDelta, 6.05) -- (12.1+\aDelta, 1.35) -- (\dCoord, 1.35);
    \draw[uniblue, very thick] (12.75+\aDelta, 6.05) -- (12.75+\aDelta, 1.35) -- (\dCoord, 1.35);   
\end{tikzpicture}


\begin{flushleft}
    \hspace{\aDelta cm}
    \begin{tabular}{|p{6.2cm}|p{2.1cm}|}
        \hline
        \rowcolor{unired}
        \multicolumn{2}{|l|}{\textcolor{white}{Модуль в слоте М12}} \\ % Объединение ячеек
        \hline
        \makecell[l]{Модуль измерительный \\ a - количество фазных ТТ с номинальным током 5А; \\ b - количество фазных ТТ с номинальным током 1А;}
          & M040.ab00 \\ 
        \hline
    \end{tabular}
\end{flushleft}

% Начинаем TikZ
\begin{tikzpicture}[overlay]
    \draw[unired, very thick] (16.75+\aDelta, 7.7) -- (16.75+\aDelta, 1.35) -- (\dCoord, 1.35); % Новые координаты
\end{tikzpicture}


\begin{flushleft}
    \hspace{\aDelta cm}
    \begin{tabular}{|p{6.2cm}|p{2.1cm}|}
        \hline
        \rowcolor{uniblue}
        \multicolumn{2}{|l|}{\textcolor{white}{Модуль в слоте М14}} \\ % Объединение ячеек
        \hline
        \makecell[l]{Модуль центрального процессора с интерфейсами:\\ Ethernet - 4 шт.} & C01.ар \\ 
        \hline
        \makecell[l]{Модуль центрального процессора с интерфейсами:\\ Ethernet - 4 шт., CANbus - 2 шт., RS485 - 2 шт.} & C02.ар \\ 
        \hline
        \multicolumn{2}{|p{8.3cm}|}{\makecell[l]{где, \textbf{a:} \\0 - 4 слота под установку SFP модулей; \\ 1 - 2xRJ45 100Base-TX + 2 слота под установку SFP модулей.}} \\ 
        \hline 
        \multicolumn{2}{|p{8.3cm}|}{\makecell[l]{где, \textbf{p:} \\
        0 - МЭК 61850-8-1; \\ 
        1 - МЭК 61850-8-1, МЭК 60870-5-104, Modbus-TCP;\\
        2 - МЭК 60870-5-104/103/101, Modbus-RTU/TCP (только для модуля C02);\\
        3 - МЭК 61850-8-1, МЭК 60870-5-104/103/101, Modbus-RTU/TCP \\(только для модуля C02);
        }} \\ 
        \hline                          
    \end{tabular}
\end{flushleft}

% Начинаем TikZ
\begin{tikzpicture}[overlay]
    \draw[uniblue, very thick] (18+\aDelta, 12.7) -- (18+\aDelta, 4.65) -- (\dCoord, 4.65); % Новые координаты
\end{tikzpicture}


\begin{flushleft}
    \hspace{\aDelta cm}
    \begin{tabular}{|p{6.2cm}|p{2.1cm}|}
        \hline
        \rowcolor{unired}
        \multicolumn{2}{|l|}{\textcolor{white}{Версия функциональной схемы устройства}} \\ % Объединение ячеек
        \hline
        \multicolumn{2}{|p{8.74cm}|}{\makecell[l]{Номер версии (справочно, опционально)}} \\ 
        \hline 
    \end{tabular}
\end{flushleft}

% Начинаем TikZ
\begin{tikzpicture}[overlay]
    \draw[unired, very thick] (18.65+\aDelta, 13.7) -- (18.65+\aDelta, 0.7) -- (\dCoord, 0.7); % Новые координаты
\end{tikzpicture}

\begin{flushleft}
    \hspace{\aDelta cm}
    \begin{tabular}{|p{6.2cm}|p{2.1cm}|}
        \hline
        \rowcolor{uniblue}
        \multicolumn{2}{|l|}{\textcolor{white}{Версия встроенного программного обеспечения}} \\ % Объединение ячеек
        \hline
        \multicolumn{2}{|p{8.74cm}|}{\makecell[l]{Номер версии (справочно, опционально)}} \\ 
        \hline 
    \end{tabular}
\end{flushleft}

% Начинаем TikZ
\begin{tikzpicture}[overlay]
    \draw[uniblue, very thick] (19.17+\aDelta, 14.7) -- (19.17+\aDelta, 0.6) -- (\dCoord, 0.6); % Новые координаты
\end{tikzpicture}







\clearpage %!!!!!!!!!!БЕЗ ЭТОГО ХРЕНЬ
\end{landscape}
\restoregeometry

